\input{preamble}

\section*{Moller scattering 1}

Moller scattering is the process $e^-+e^-\rightarrow e^-+e^-$.

\begin{center}
\begin{tikzpicture}
\draw[dashed] (0,0) circle (0.5cm);
\draw[thick,->] (2,0) node[anchor=west] {$e^-$} -- (0.6,0);
\draw[thick,->] (-2,0) node[anchor=east] {$e^-$} -- (-0.6,0);
\draw[thick,->] (0.40,0.40) -- (1.3,1.3) node[anchor=south west] {$e^-$};
\draw[thick,->] (-0.4,-0.4) -- (-1.3,-1.3) node[anchor=north east] {$e^-$};
\draw (1,0.5) node {$\theta$};
\end{tikzpicture}
\end{center}

The following center-of-mass momentum vectors have $E=\sqrt{p^2+m^2}$.
\begin{equation*}
p_1=\underset{e^- \, \longrightarrow}
{\begin{pmatrix}E\\0\\0\\p\end{pmatrix}}
\qquad
p_2=\underset{\longleftarrow \, e^-}
{\begin{pmatrix}E\\0\\0\\-p\end{pmatrix}}
\qquad
p_3=\underset{\substack{\phantom{e^-} \, \nearrow\\e^- \, \phantom{\nearrow}}}
{\begin{pmatrix}
E\\
p\sin\theta\cos\phi\\
p\sin\theta\sin\phi\\
p\cos\theta
\end{pmatrix}}
\qquad
p_4=\underset{\substack{\phantom{\swarrow} \, e^-\\\swarrow \, \phantom{e^-}}}
{\begin{pmatrix}
E\\
-p\sin\theta\cos\phi\\
-p\sin\theta\sin\phi\\
-p\cos\theta
\end{pmatrix}}
\end{equation*}

Spinors for $p_1$.
\begin{equation*}
u_{11}=\frac{1}{\sqrt{E+m}}
\underset{\text{spin up }}
{\begin{pmatrix}E+m\\0\\p\\0\end{pmatrix}}
\qquad
u_{12}=\frac{1}{\sqrt{E+m}}
\underset{\text{spin down }}
{\begin{pmatrix}0\\E+m\\0\\-p\end{pmatrix}}
\end{equation*}

Spinors for $p_2$.
\begin{equation*}
u_{21}=\frac{1}{\sqrt{E+m}}
\underset{\text{spin up}}
{\begin{pmatrix}E+m\\0\\-p\\0\end{pmatrix}}
\qquad
u_{22}=\frac{1}{\sqrt{E+m}}
\underset{\text{spin down}}
{\begin{pmatrix}0\\E+m\\0\\p\end{pmatrix}}
\end{equation*}

Spinors for $p_3$.
\begin{equation*}
u_{31}=\frac{1}{\sqrt{E+m}}
\underset{\text{spin up}}
{\begin{pmatrix}E+m\\0\\p_{3z}\\p_{3x}+ip_{3y}\end{pmatrix}}
\qquad
u_{32}=\frac{1}{\sqrt{E+m}}
\underset{\text{spin down}}
{\begin{pmatrix}0\\E+m\\p_{3x}-ip_{3y}\\-p_{3z}\end{pmatrix}}
\end{equation*}

Spinors for $p_4$.
\begin{equation*}
u_{41}=\frac{1}{\sqrt{E+m}}
\underset{\text{spin up}}
{\begin{pmatrix}E+m\\0\\p_{4z}\\p_{4x}+ip_{4y}\end{pmatrix}}
\qquad
u_{42}=\frac{1}{\sqrt{E+m}}
\underset{\text{spin down}}
{\begin{pmatrix}0\\E+m\\p_{4x}-ip_{4y}\\-p_{4z}\end{pmatrix}}
\end{equation*}

The scattering amplitude $\mathcal M$ is
\begin{equation*}
\mathcal M=\mathcal M_1-\mathcal M_2
\end{equation*}

where $\mathcal M_1$ is the amplitude for the $1\rightarrow3$, $2\rightarrow4$ process
\begin{equation*}
\mathcal M_1=\frac{\alpha}{t}
(\bar{u}_{4d}\gamma^\mu u_{2b})
(\bar{u}_{3c}\gamma_\mu u_{1a})
\end{equation*}

and $\mathcal M_2$ is the amplitude for the $1\rightarrow4$, $2\rightarrow3$ process
\begin{equation*}
\mathcal M_2=\frac{\alpha}{u}
(\bar{u}_{3c}\gamma^\nu u_{2b})
(\bar{u}_{4d}\gamma_\nu u_{1a})
\end{equation*}

The expected density $\langle|\mathcal M|^2\rangle$ is the average over spin states.
\begin{equation*}
\langle|\mathcal M|^2\rangle=\frac{1}{4}\sum_{abcd}|\mathcal M|^2
\end{equation*}

Expand the summand.
\begin{equation*}
\langle|\mathcal{M}|^2\rangle=\frac{1}{4}
\sum_{abcd}
\left(
\mathcal M_1\mathcal M_1^*-
\mathcal M_1\mathcal M_2^*-
\mathcal M_2\mathcal M_1^*+
\mathcal M_2\mathcal M_2^*
\right)
\end{equation*}

The summation results are
\begin{align*}
\sum_{abcd}
\mathcal M_1\mathcal M_1^*&=\frac{\alpha^2}{t^2}
\left(8 s^2 + 8 u^2 - 64 m^2 s - 64 m^2 u + 192 m^4\right)
\\
\sum_{abcd}
\mathcal M_1\mathcal M_2^*&=\frac{\alpha^2}{tu}
\left(-8 s^2 + 64 m^2 s - 96 m^4\right)
\\
\sum_{abcd}
\mathcal M_2\mathcal M_2^*&=\frac{\alpha^2}{u^2}
\left(8 s^2 + 8 t^2 - 64 m^2 s - 64 m^2 t + 192 m^4\right)
\end{align*}

where $s$, $t$, and $u$ are the Mandelstam variables
\begin{align*}
s&=(p_1+p_2)^2
\\
t&=(p_1-p_3)^2
\\
u&=(p_1-p_4)^2
\end{align*}

For clarity define
\begin{align*}
f_{11}&=8 s^2 + 8 u^2 - 64 m^2 s - 64 m^2 u + 192 m^4
\\
f_{12}&=-8 s^2 + 64 m^2 s - 96 m^4
\\
f_{22}&=8 s^2 + 8 t^2 - 64 m^2 s - 64 m^2 t + 192 m^4
\end{align*}

Now write $\langle|\mathcal{M}|^2\rangle$ more compactly as
\begin{equation*}
\langle|\mathcal{M}|^2\rangle
=\frac{\alpha^2}{4}
\left(\frac{f_{11}}{t^2}-\frac{2f_{12}}{tu}+\frac{f_{22}}{u^2}\right)
\end{equation*}

For $E\gg m$ a useful approximation is to set $m=0$ and obtain
\begin{align*}
f_{11}&=8s^2+8u^2\\
f_{12}&=-8s^2\\
f_{22}&=8s^2+8t^2
\end{align*}

Then $\langle|\mathcal{M}|^2\rangle$ simplifies as
\begin{equation*}
\langle|\mathcal{M}|^2\rangle
=2\alpha^2
\left(\frac{s^2+u^2}{t^2}+\frac{2s^2}{tu}+\frac{s^2+t^2}{u^2}\right)
\end{equation*}

For $m=0$ the Mandelstam variables are
\begin{align*}
s&=4E^2
\\
t&=-2E^2(1-\cos\theta)
\\
u&=-2E^2(1+\cos\theta)
\end{align*}

Hence
\begin{equation*}
\langle|\mathcal{M}|^2\rangle
=2\alpha^2\left(
\frac{1+\cos^4(\theta/2)}{\sin^4(\theta/2)}
+\frac{2}{\sin^2(\theta/2)\cos^2(\theta/2)}
+\frac{1+\sin^4(\theta/2)}{\cos^4(\theta/2)}
\right)
\end{equation*}

This can be written more compactly as
\begin{equation*}
\langle|\mathcal{M}|^2\rangle=4\alpha^2\frac{(\cos^2\theta+3)^2}{\sin^4\theta}
\end{equation*}

The differential cross section is
\begin{equation*}
\frac{d\sigma}{d\Omega}=\frac{\langle|\mathcal{M}|^2\rangle}{4s}
=\frac{\alpha^2}{4E^2}\frac{(\cos^2\theta+3)^2}{\sin^4\theta}
\end{equation*}

Multiply by $(\hbar c)^2$ to convert from natural units to physical units.
\begin{equation*}
\frac{d\sigma}{d\Omega}=\frac{\alpha^2}{4E^2}
\frac{(\cos^2\theta+3)^2}{\sin^4\theta}(\hbar c)^2
\end{equation*}

Zee p.~134 has the cross section formula
\begin{equation*}
\frac{d\sigma}{d\Omega}=\left(\frac{e^2}{4\pi}\right)^2\frac{1}{8E^2}f(\theta)
\end{equation*}
where $f(\theta)$ is the function
\begin{equation*}
f(\theta)=
\frac{1+\cos^4(\theta/2)}{\sin^4(\theta/2)}
+\frac{2}{\sin^2(\theta/2)\cos^2(\theta/2)}
+\frac{1+\sin^4(\theta/2)}{\cos^4(\theta/2)}
\end{equation*}

Function $f(\theta)$ is equivalent to
\begin{equation*}
f(\theta)=2\,\frac{(\cos^2\theta+3)^2}{\sin^4\theta}
\end{equation*}

Noting that in natural units $e^2=4\pi\alpha$ we have
\begin{equation*}
\left(\frac{e^2}{4\pi}\right)^2=\alpha^2
\end{equation*}

Hence
\begin{equation*}
\frac{d\sigma}{d\Omega}=\frac{\alpha^2}{4E^2}
\frac{(\cos^2\theta+3)^2}{\sin^4\theta}
\end{equation*}

\href{https://georgeweigt.github.io/examples/moller-scattering-1-demo.html}{Eigenmath script}

\end{document}
